\section{The Always-On Evaluation Protocol (AOEP)}\label{sec:aoep}


\regime{Control}{how governance can be scored on a running system rather than only described.}
The preceding parts established a recurring asymmetry in our coded corpus. The literature has learned to accumulate state and to retrieve it, and a large benchmark literature now scores how well it does so, but much less of that literature scores whether the state was \emph{governed}: whether a deleted fact stayed deleted across derived tiers, whether an action was licensed by a still-valid authority, whether an untrusted observation was prevented from becoming durable instruction, whether a corrupted decision can be traced back and undone. The Always-On Evaluation Protocol (AOEP) targets that omission. It does not measure answer quality on a memory question; it measures the correctness of a \emph{state trajectory} under the kinds of perturbations only a persistent agent encounters: restart, conflict, deletion request, permission revocation, shared-scope change, adversarial write, and delayed consequence. AOEP-v0 is a reference design plus a pilot analysis, with the evaluation flow shown in Figure~\ref{fig:aoep}. The pilot illustrates how a survey-derived contract can expose missing governance state in tested memory configurations; it is not evidence that deployed memory systems generally fail or that larger readers could never help once the memory layer exposes the right fields.

AOEP is \emph{representational, not a leaderboard}, and a \emph{wrapper around realistic substrates, not a replacement for them}. AOEP-v0 is not intended to rank deployed memory products or establish benchmark-grade generalization. The contribution is a contract for what a state transition must expose to be scorable, plus a small harness showing how that contract can distinguish stores that expose governance fields from stores that expose only text.

AOEP operationalizes the six axes by giving each one concrete fields in an event stream and state snapshot. Authority becomes permission epochs and approval records; scope becomes actor, resource, privacy class, and sharing fields; mutability becomes typed update, supersession, conflict, and tombstone operations; provenance becomes causal links, trust tiers, signatures, and source identifiers; recoverability becomes idempotency keys, transaction ids, deletion ledgers, and rollback ledgers; and actionability becomes the typed operation attached to the record or side effect. The protocol therefore does not add a separate taxonomy. It turns the earlier axes and lifecycle invariants into the fields a validator can check.

\subsection{Why existing evaluation cannot score governance}

The evaluation literature has strong machinery for two questions and much less for a third. On judge reliability, a long line of work establishes the statistical machinery for trustworthy comparison: \citet{zheng2023judgingllmasajudge} establishes LLM-as-judge via MT-Bench and Chatbot Arena, \citet{gu2024surveyllmasajudge} and \citet{li2024llmsasjudgessurvey} synthesize consistency and bias-mitigation methodology, and \citet{shi2024judgingthejudges} quantifies position bias with repetition-stability and position-consistency diagnostics. \citet{you2026agentasjudge} extends the paradigm from judging outputs to judging agentic trajectories step-by-step. This body of work tells us how to score reliably; it does not tell us \emph{what} state property to score. On the methodology of agent evaluation itself, \citet{kapoor2024aiagentsthatmatter} exposes the absence of cost-versus-accuracy reporting and the overfitting that follows from inadequate holdout sets, \citet{miller2024errorbarsevals} and \citet{madaan2024quantifyingvariance} supply confidence intervals, two-model difference tests, and empirical seed-variance estimates, and \citet{white2024livebench} addresses contamination through monthly-refreshed questions. These works discipline how we report; they assume the unit being reported is task success.

On memory specifically, a 2026 wave has begun to pull retrieval correctness apart from answer correctness. \citet{flynt2026precisionmembench} argues that LoCoMo-style benchmarks conflate answer-correctness with retrieval-correctness and contributes PrecisionMemBench to score retrieval in isolation from the generator. \citet{long2026memtrace} re-bases long-term-memory evaluation on the per-fact knowledge point probed across age, question-type, and evidence conditions, showing that pooled accuracy hides distinct failure modes and that the bottleneck is evidence use, not retrieval. \citet{liu2026streammembench} diagnoses, over egocentric streams, whether agents turn stored evidence and feedback into reliable future behavior, and \citet{shao2026scaleconditioned} holds task evidence fixed while adding irrelevant sessions to find where stored evidence stops being usable. This is real progress, but its frame remains \emph{retrieval and use of evidence for an answer}. Compared against each other, these benchmarks differ in what they isolate (retrieval precision, per-fact tracing, streaming feedback, distractor scaling) yet they share a blind spot: none records whether the value that was retrieved was authorized to act, whether its source was trusted, or whether a later deletion request was honored across the store.

Two recent works move closest to the state-trajectory view and frame the gap precisely. \citet{jung2026meme} contributes a multi-entity evolving-memory benchmark whose Cascade, Absence, and Deletion tasks test whether updates and removals propagate to dependent entries, and reports near-total collapse (Cascade 3\%, Absence 1\%) of current systems on dependency reasoning. \citet{kwon2026reclaim} shows, judge-free, that a memory keeping a stale conclusion but dropping its source becomes confidently uncorrectable and strictly worse than empty memory, fixed only by a source-first write policy. Most directly, \citet{orogat2026gem} reframes long-term agent-memory correctness as a property of the state trajectory rather than individual records, formalizing Governed Evolving Memory with state-level operators and six correctness conditions, and proving that record-level stores cannot satisfy them. AOEP should be read as a prototype instantiation of this state-trajectory view: where \citet{orogat2026gem} supplies a formal impossibility result, AOEP supplies a checkable contract that can be wrapped around concrete memory systems, and where \citet{jung2026meme} measures deletion propagation alone, AOEP scores it alongside authority currency, provenance trust, conflict surfacing, and rollback logging. The gap this subsection exposes is concrete: state-trajectory correctness has been formalized and partially benchmarked, but we found no widely adopted protocol that drives a memory system through restart, conflict, deletion, revocation, and adversarial write while scoring all five governance invariants deterministically. That is the role AOEP-v0 is meant to prototype.

\subsection{Protocol: event stream and state snapshot}

AOEP fixes a data contract rather than a task. An episode is a sequence of typed \emph{events}, replayed to a system under test one at a time, and a \emph{state snapshot} reconstructed at the end from the system's answers to neutral probes. The protocol's claim is that an event must carry strictly more than a memory dump records, and that this added structure is what makes governance scorable.

Concretely, each AOEP event carries fields a flat memory entry omits. It carries an \emph{idempotency key}, so a replayed duplicate write after a restart is recognizable as the same logical operation rather than a second fact. It carries \emph{causal links}, separated into parent, supersedes, and conflicts-with edges plus a transaction id, so that a later value can be marked as superseding an earlier one and a collaborator's proposal can be marked as conflicting with an owner's. It carries a \emph{permission epoch} on its scope, an integer that increments when authority changes, so that an action can be checked against the epoch in force at execution rather than the epoch at write. It carries \emph{provenance} with an explicit \emph{trust tier} and an optional signature, so an untrusted tool result is distinguishable from a user instruction regardless of how persuasive its text is. It carries \emph{retention and privacy constraints}, a retention policy, a privacy class, an optional time-to-live, and a requires-confirmation flag, so that deletion and confirmation obligations attach to the record itself. And it carries one \emph{typed operation} drawn from a closed set of eleven: \texttt{read}, \texttt{write}, \texttt{update}, \texttt{delete}, \texttt{tombstone}, \texttt{share}, \texttt{unshare}, \texttt{validate}, \texttt{quarantine}, \texttt{deny}, and \texttt{rollback}. The closed operation set is what lets the validator reason about the trajectory as a state machine: a \texttt{tombstone} that is not reflected in the deletion ledger, or a \texttt{quarantine} that is later promoted to instruction, is a named, locatable violation rather than a degraded score.

The snapshot is the dual of the stream. It exposes a \emph{deletion ledger}, a \emph{rollback ledger}, the set of \emph{pending conflicts} with their candidate values and the authority backing each, the \emph{per-resource permission epochs} currently in force, and the boolean results of the invariant checks. The protocol deliberately separates the system's self-report from the harness's verdict: the harness recomputes every invariant boolean itself from the reconstructed snapshot and never trusts a system's claim that it behaved. This design borrows the spirit of trajectory-level evaluation from \citet{you2026agentasjudge} while deliberately avoiding LLM-as-judge scoring, because the position, verbosity, and self-enhancement biases documented by \citet{zheng2023judgingllmasajudge} and \citet{shi2024judgingthejudges} would be fatal for a contract whose entire purpose is to detect quiet governance failures. Every AOEP check is deterministic.

This design choice distinguishes a state snapshot from a memory dump, and it directly answers the worked example threaded through AOEP. An episode begins with a user setting a durable billing value; the system restarts and replays a duplicate write; a collaborator proposes a conflicting value in shared scope; the owner then requests deletion. A memory-QA scorer marks the trajectory correct once the final recalled value is right. AOEP fails it if the tombstoned value survives in a derived tier (deletion propagation) or the write was licensed by a lapsed authority epoch (authority monotonicity): distinctions a recall unit has no field in which to record. The gap is tied to the thesis: the wire formats the field actually ships do not carry these fields. The vendor-neutral memory wire format of the substrate literature defines remember/recall/forget/merge/expire over text but specifies neither permission epoch, trust tier, nor causal supersession, so even an interoperable memory layer has nowhere to put the state on which governance depends.

\subsection{Deterministic coverage scorecard}

AOEP's scorecard is deliberately split into two scores rather than a single number. The first is \emph{obligation pass}: the positive obligations a system must \emph{actively} satisfy. These are operations that produce state the agent must maintain, namely recording a deletion, reporting the current permission epoch after a revocation, blocking a stale-permission or untrusted-instruction task, surfacing an owner-versus-collaborator conflict, and logging a rollback after an external action. The second is \emph{negative-invariant pass}: the no-leakage checks, namely that a deleted value is not visible, that an out-of-scope value does not appear, and that an untrusted instruction was not promoted.

The reason the split matters is that the two scores have opposite degenerate solutions. A system that stores nothing trivially passes every negative invariant: it cannot leak a deleted address because it never retained one, and it cannot promote an untrusted instruction because it never kept any text. If a single scalar mixed positive and negative checks, a no-memory floor would look respectable, and the protocol would reward amnesia. By separating them, AOEP makes the no-memory floor legible as what it is: a system that never leaks precisely because it satisfies no positive obligation at all. The headline must therefore be obligation pass, with negative-invariant pass read as a guardrail that catches the opposite failure (a system that recalls so eagerly it leaks). This mirrors, at the protocol level, the lesson of \citet{kwon2026reclaim} that empty memory can dominate confidently-wrong memory, and the lesson of \citet{flynt2026precisionmembench} that a pooled score hides which side of the tradeoff a system actually failed. The validator maps the five named invariants of the lifecycle to ten executable checks (eight per-snapshot booleans plus two ledger-subset checks): authority monotonicity decomposes into \texttt{no\_unauthorized\_writes}, \texttt{permission\_epoch\_current}, and \texttt{no\_stale\_action\_executed}; scope non-expansion into \texttt{no\_scope\_leakage}; deletion propagation into \texttt{no\_deleted\_content\_visible} plus a deletion-ledger subset match; provenance preservation into \texttt{no\_untrusted\_instruction\_promoted}; and rollback traceability into a rollback-ledger subset match plus \texttt{no\_external\_action\_without\_approval}. A non-conforming trajectory fails a specific named check, not a scalar, which is what makes the scorecard diagnostic rather than only comparative.

\subsection{Pilot harness}

To test whether the protocol discriminates real systems rather than only validating hand-authored reference trajectories, we built a harness around the same schema and ran seven systems through nine distinct fault patterns: the three hand-authored worked episodes (restart-conflict-deletion, permission-epoch drift, adversarial memory injection) plus six new patterns that each stress a different governance situation (deletion-derived summary residue, cross-user scope leak, stale permission after a restart, owner/collaborator conflict resolution, rollback of an external action, and untrusted tool-output poisoning). The harness strips outcome-revealing fields from each event, delivers the stream one event at a time, lets each system maintain its state however it chooses, and at the end poses neutral probes that name only target and actor identities, never expected values or invariant names. It then reconstructs the snapshot, computes the invariants itself, and scores against the validator.

Seven systems run. A deterministic \emph{governed reducer} applies AOEP rules directly; it is an oracle-free upper bound, not a deployable product. A \emph{no-memory floor} answers from nothing. Five ungoverned memory systems run over a local Qwen2.5-7B reader: \emph{naive append} (store all event text, retrieve by recency), \emph{full context} (all events in the prompt), \emph{vector-RAG} (all-MiniLM embeddings, top-$k$ cosine), a \emph{Mem0-style} reimplementation in which the reader extracts atomic memory facts from each event into a vector index, and \emph{Mem0 (package)}, a local \texttt{mem0ai} configuration run through a local OpenAI-compatible endpoint backed by the same reader, a local embedder, and an in-process vector store. All computation is local and decoding is greedy, so runs are deterministic. Including this local \texttt{mem0ai} setup reduces the narrow objection that the reimplementation alone created the failure mode; it does not turn the pilot into a product evaluation. The connection to production memory layers is mainly methodological: cost studies such as \citet{wolff2026costaccuracy} evaluate similar store families without pricing mutation or deletion-correctness obligations.

\begin{table}[t]
\centering
\small
\begin{tabular}{lccc}
\toprule
System & Obligation & Neg.-invariant & All \\
\midrule
Governed reducer (upper bound) & $15/15$ & $41/41$ & $92/92$ \\
No-memory (floor) & $0/15$ & $41/41$ & $77/92$ \\
Naive append (Qwen2.5-7B) & $7/15$ & $40/41$ & $83/92$ \\
Full context (Qwen2.5-7B) & $7/15$ & $40/41$ & $83/92$ \\
Vector-RAG (MiniLM+Qwen2.5-7B) & $7/15$ & $40/41$ & $83/92$ \\
Mem0-style (extracted facts) & $4/15$ & $38/41$ & $77/92$ \\
Mem0 (actual \texttt{mem0ai}) & $3/15$ & $36/41$ & $73/92$ \\
\bottomrule
\end{tabular}
\caption{AOEP pilot over nine fault patterns, one frozen reader, greedy decoding. The governed upper bound satisfies every obligation. The three raw-storage systems satisfy the same recall obligations and fail the same governance obligations, scoring identically across recency, full-context, and dense retrieval. The tested extracted-fact configurations score lower (Mem0-style $4/15$; the local \texttt{mem0ai} setup $3/15$), because in this pilot their extraction step drops the structured envelope. The no-memory floor is the only system that never leaks, and it does so by satisfying no positive obligation at all.}
\label{tab:aoep-pilot}
\end{table}

Table~\ref{tab:aoep-pilot} reports the headline. The reading turns entirely on the obligation column. The governed reducer satisfies all fifteen obligations, demonstrating that the obligations are satisfiable inside the reference harness; the no-memory floor satisfies none. The three raw-storage systems all score $7/15$, and the identity of their scores is the most informative fact in the table. They pass the obligations that reduce to \emph{semantic recall}, namely whether a value was deleted, whether an instruction was quarantined, and the status of an active item, all of which a capable reader infers from text. They fail the obligations that require maintained \emph{governance state}: reporting the current permission epoch after a revocation, blocking a stale-permission or untrusted-instruction task, surfacing a conflict, and logging a rollback. None of these is a retrieval problem in the tested configurations, which is why recency, full-context, and dense retrieval score identically. The information is not missing from the context; it is that these text stores have no field in which a permission epoch, a quarantine state, a conflict record, or a rollback ledger lives. This is a pilot instantiation of the state-trajectory concern formalized by \citet{orogat2026gem}: three stores that differ in retrieval strategy fail the same governance subset when the required envelope is absent.

The extracted-fact systems sharpen the point by scoring \emph{below} raw storage in this pilot. The Mem0-style store reaches $4/15$ and the local \texttt{mem0ai} configuration reaches $3/15$, because extracting atomic facts into a vector index flattens the structured event (its permission epoch, deletion link, and trust tier) into free text, so even the deletion-awareness that raw storage incidentally keeps can be lost. In these tested configurations, extraction without an explicit governance envelope performs worse on governance obligations, not better. The local package setup leaks a deleted billing address and an untrusted exfiltration address into its retrieved memories on two episodes, which is the cross-session injection threat and the deletion-residue failure made measurable in a concrete implementation. This is not a defect claim about \texttt{mem0ai} as software, nor a verdict on extracted-fact architectures in principle. It is a scoped demonstration of the failure that source-first write policies \citep{kwon2026reclaim} and dependency-propagation tasks \citep{jung2026meme} would predict when the envelope is absent.

\subsection{Robustness of the pilot}

The pilot result invites three objections, each of which would make the observed gap an artifact of setup rather than a property of the tested designs.

\emph{Reader variation does not remove the pilot gap.} Running the full-context system, where the entire event stream sits in the prompt so any failure is missing governance state rather than missing retrieval, across four readers spanning sizes and families holds obligation pass at or below the original while the largest tested reader does not help. Table~\ref{tab:aoep-reader} reports the span: Qwen2.5-3B, Qwen2.5-7B, and Llama-3.1-8B all score $7/15$, Mistral-7B-Instruct scores $4/15$, and the governed reducer scores $15/15$ on the same patterns. Within this local sweep, the missing objects are not facts hidden from the reader but explicit fields the text store never maintained: an epoch counter, a deletion ledger, and a conflict record. This is a pilot-level version of the separation \citet{long2026memtrace} draws when it finds the bottleneck to be evidence use rather than retrieval, and the same one \citet{shao2026scaleconditioned} draws when stored evidence stops being usable independent of model strength.

\begin{table}[t]
\centering
\small
\begin{tabular}{lcc}
\toprule
Full-context reader & Obligation & Neg.-invariant \\
\midrule
Qwen2.5-3B-Instruct & $7/15$ & $34/41$ \\
Qwen2.5-7B-Instruct & $7/15$ & $40/41$ \\
Mistral-7B-Instruct-v0.3 & $4/15$ & $39/41$ \\
Llama-3.1-8B-Instruct & $7/15$ & $38/41$ \\
\midrule
Governed reducer (upper bound) & $15/15$ & $41/41$ \\
\bottomrule
\end{tabular}
\caption{Reader robustness over the nine fault patterns, full-context system. In this local $3$B to $8$B sweep across three model families, no tested reader approaches the governed upper bound on governance obligations; the larger reader does no better and the smaller one leaks more. The result is consistent with a missing-governance-state explanation, but it is not a model-scaling claim.}
\label{tab:aoep-reader}
\end{table}

\emph{It is not a single-run effect.} The main runs decode greedily. Re-running the full-context system under sampled decoding (temperature $0.7$, top-$p$ $0.95$) across five seeds holds obligation pass in a $6$ to $8/15$ band (mean $6.6$, standard deviation $0.8$), straddling the greedy $7/15$ and nowhere near the governed $15/15$. The governance gap is stable to sampling variance, which is the kind of seed-sensitivity control that \citet{madaan2024quantifyingvariance} argues every agent result should report and that \citet{miller2024errorbarsevals} formalizes; we report it here so the separation is not a single lucky draw.

\emph{It is not an artifact of synthetic fault patterns.} To check that the split survives realistic task shapes, we authored three longer, multi-actor episodes from documented agent-task benchmarks: a tau-bench-style customer-refund flow (approval policy, an untrusted chat message redirecting the refund, a stale retry, a card-deletion request), a TheAgentCompany-style enterprise-scheduling flow (granted then revoked calendar access, a contested reschedule, a stale write), and one grounded in an actual AppWorld task \citep{trivedi-etal-2024-appworld}, the group-dinner Venmo payment-request task, using AppWorld's real apps and operations with an untrusted note redirecting the payment and a card-deletion request. The pattern holds and sharpens: across the three episodes the governed reducer satisfies all nine obligations, the three raw-storage systems satisfy four and the Mem0-style store two, and the raw-storage systems leak the attacker's redirect instruction (and, on the Venmo task, the stored payment card) into the action flow. The recall-versus-governance gap is a property of the memory designs, not of the minimal fault patterns. This matters because \citet{chen2026crossscenario} shows that memory systems tuned to one trajectory format collapse across five others; the persistence of the AOEP split across hand-authored, realistic, and benchmark-grounded shapes gives the same warning from the governance side.

\emph{It reproduces on real execution traces.} The realistic episodes are hand-authored from real tasks; to ground the base streams in real executions we recorded genuine AppWorld traces by driving its actual API server across three apps with verified side effects. On the Venmo task, \texttt{create\_payment\_request} returned a live request id and \texttt{delete\_payment\_card} dropped the account from four stored cards to three; on a Spotify library-cleanup task, \texttt{remove\_song\_from\_library} dropped the library from fourteen songs to thirteen and the song stayed absent; on a file-system task, a create-then-delete created and removed a file holding a secret, verified present then gone. Each real call was captured as an AOEP event, and an untrusted note was added (redirect the payment, re-add the deleted song, recreate the deleted file). Across the three traces the governed reducer satisfies all seven obligations while the ungoverned systems satisfy three to six, and each leaks an untrusted instruction on one trace. The base operations actually executed against the API backend with verified side effects; only the perturbations and probes are authored. Finally, on \emph{agent-chosen} trajectories, a stronger local policy (Qwen2.5-Coder-32B) drove AppWorld autonomously, fetching credentials, logging in, and completing three single-app tasks in its own self-chosen steps (a harder multi-step task it did not finish, which we report rather than hide); when one untrusted observation was injected mid-trajectory, the governed reducer handled it on all three completed runs while the full-context system leaked the exfiltration instruction. The gap appears even on steps the agent selected rather than ones we scripted.

\subsection{What the pilot establishes, and what it leaves open}

Read together, the pilot and its robustness checks support a scoped qualitative claim: recall and governance separate once the evaluation contract asks for fields such as permission epochs, deletion ledgers, trust tiers, conflict records, and rollback ledgers. Recall is necessary and the raw-storage systems have it; governance requires representing state that a flat or extracted store may not carry unless it is explicitly instrumented to do so. In this local sweep, no tested reader from 3B to 8B across three families manufactures that state from text it can read. The local \texttt{mem0ai} configuration failing below raw storage makes the point concrete, but the paper should not be read as ranking that package or as establishing class-wide performance of extracted-fact systems.

\paragraph{A required ablation before stronger claims.} The pilot should be read as a configuration-level result unless and until one ablation is run: add a minimal governance envelope to an extracted-fact store, consisting of a permission epoch, deletion ledger, and provenance trust tier, and rerun the same nine patterns. This is the unit test that separates two different claims. If the envelope closes most of the gap, then the result is that current tested configurations omit governance metadata. If it does not, then extraction itself may be discarding structure in a way that the envelope cannot repair. The present manuscript supports only the first, weaker reading: in the tested configurations, the fields required by AOEP are absent, so readers and retrieval strategies cannot recover them.

The limitations follow from that gate. This is a pilot, not a benchmark: nine fault patterns plus three realistic episodes and three real execution traces, scored over a handful of readers and memory systems. AOEP is built to wrap realistic substrates, and the natural next step is to drive it across hosted, multi-task agent benchmarks at the scale that \citet{kapoor2024aiagentsthatmatter} demands for cost-aware reporting and that the snapshot-collapse phenomenon of \citet{zheng2026seagym} warns must be measured over long horizons rather than single snapshots. The pilot measures whether ungoverned systems \emph{do} expose the required governance state, not whether they \emph{could} govern if minimally augmented. It scores one local \texttt{mem0ai} configuration as run, and that configuration carries no permission epoch, deletion ledger, or trust tier. We did not run the envelope ablation above, so we do not claim to separate an architectural limit from an instrumentation gap. The deeper gap, tied directly to the thesis, is the one the corpus quantifies as a scoping estimate: rollback is studied as a mechanism but rarely measured, with only $27$ of $435$ works exposing any rollback mechanism and none in our corpus reporting recovery success or cost after corruption. AOEP scores rollback as one obligation among fifteen; what it does not yet do is measure the \emph{cost} of recovery, the latency and side-effect-replay correctness of undoing a corrupted decision once the authority that licensed it has lapsed. Prospective and spontaneous governance is a second open axis: \citet{zhang2026triggerbench} shows that spontaneous recall and action on a latent stored constraint is far harder than answering a direct query and degrades with context length. An always-on agent therefore has to honor deletion or revocation when asked and also surface the obligation unprompted when the context demands it, a capability AOEP can perturb through its delayed-consequence pattern but does not yet score at scale. The evaluation agenda is to measure recovery cost and prospective governance over realistic substrates with the statistical rigor of \citet{miller2024errorbarsevals} and \citet{madaan2024quantifyingvariance}. AOEP's specific advance over the survey's earlier parts is that it turns the diagnosis into a checkable contract for what a state transition must expose.
